\section{Introduction}
\label{sec:introduction}

AIMer is a post-quantum signature scheme based on a symmetric one-way function
and an MPC-in-the-Head (MPCitH) proof.  Its third specification replaces AIM2
with the new AIM3 primitive.  AIM3 introduces an IV-dependent input affine
layer, enhanced Mersenne S-boxes, an output affine layer, and zero-input
rejection during key generation~\cite{cho2026aimer3}.  These changes respond to
new algebraic analyses of AIM2, including the Magic Pot
attack~\cite{biryukov2026magic} and subsequent studies of weak AIM2 parameter
choices~\cite{sun2025high}.

The transition is substantial from a cryptographic-design perspective, but it
does not remove the implementation-level parallelism of AIMer.  Signing and
verification still simulate many independent MPC parties.  Each party performs
the same operations over $\gftwo{128}$, $\gftwo{192}$, or $\gftwo{256}$, and
commitment and tape generation still invoke many independent SHAKE instances.
These common properties suggest that the \emph{optimization principles} used
for AIMer v2 can be transferred even though its algorithm-specific code cannot
be reused unchanged.

This distinction is important.  Reusing AIM2 constants, exponentiation chains,
or MPC equations would produce a different algorithm and invalidate the AIMer
v3 test vectors.  Our approach instead preserves the AIM3 computation and
changes only its evaluation schedule and register layout.  We implement AVX2
and AVX-512 over the same high-level scope so that their performance can be
compared directly.

The contributions of this paper are as follows.

\begin{enumerate}[label=\arabic*.,leftmargin=*]
\item We present matched-scope AVX2 and AVX-512 implementations of all six
AIMer v3 parameter sets and expose them through a self-contained
OQS-compatible library with run-time dispatch.

\item We transfer four implementation principles from earlier AIMer work:
four-way SHAKE, party-parallel field arithmetic, branch-free binary
matrix--vector multiplication, and operation-wise MPC batching.  We describe
the concrete register layouts and instruction choices for both instruction
sets.

\item We establish implementation equivalence through 1,800 byte-exact KAT
responses, differential tests of scalar and batched field operations, API and
tamper tests, and sanitizer builds.

\item We evaluate key generation, signing, and verification under a fixed and
pinned single-core environment.  AVX-512 is $1.43$--$1.52\times$ faster than
AVX2 for signing and verification while both substantially outperform the
portable reference implementation.
\end{enumerate}

To the best of our knowledge, as of September 2026 no prior publication
reports a matched AVX2/AVX-512 implementation and end-to-end evaluation of
AIMer v3.  We do not claim the first SIMD implementation of AIMer, binary-field
arithmetic, or Keccak; our contribution is their careful transfer to the new
AIM3 computation and a controlled comparison over an identical scope.

The remainder of this paper is organized as follows.  Section~\ref{sec:bg}
summarizes AIMer v3 and related work.  Section~\ref{sec:impl} describes the
implementation.  Section~\ref{sec:eval} presents correctness and performance
results, and Section~\ref{sec:conc} concludes the paper.
